\begin{pastexam}[label={px:df3-098-ds-03}]{地方三等資訊處理 98 年 第三題}
（申論；25 分）

二元搜尋樹（binary search tree）：
\begin{enumerate}[leftmargin=2.2em]
\item[㈠] 二元搜尋樹與堆積（heap）有何主要相異點？寫出二項。（6 分）
\item[㈡] 將一含有 n 個節點（n>1）之二元搜尋樹以堆積來表示，並以一陣列來儲存此堆積，請問此陣列容量
可能之最小值與最大值分別為何？請說明原因。（6 分）
\item[㈢] 已知二元樹節點含有二個指標欄位以指向其左子與右子，請問一棵具有 n 個節點（n>=1）的二元樹，
存在多少個空的指標欄位（也就是其值為 NULL）？請說明原因。（5 分）
\item[㈣] 依序將整數鍵值 45, 33, 17, 65, 54, 70, 88, 25 加入一棵空的二元搜尋樹，再繪製此二元樹對應之
引線樹（threaded tree）。（8 分）
\end{enumerate}
\end{pastexam}

\begin{pastsoln}
㈠ 屬於資結書第~8、9~章，㈡ ㈢ ㈣ 是二元樹的表示法和引線樹（第~7~章），所以歸在第~7~章。

\textbf{㈠ 相異點（寫兩項即可）。} 以最大堆積為例：
\begin{itemize}
\item \textbf{排序性質不同。} 二元搜尋樹：左子樹的鍵都比根小、右子樹的都比根大，中序走訪就是
由小到大。 堆積：只要求父節點不小於子節點，左右兩個子節點之間沒有大小關係，中序走訪沒有意義。
\item \textbf{形狀限制不同。} 堆積一定是完整二元樹（complete binary tree），高度固定是
$\lfloor \lg n \rfloor + 1$ 層，可以用沒有空洞的陣列存。 二元搜尋樹的形狀由插入順序決定，可能很平衡
也可能歪成一條鏈，高度在 $\lfloor \lg n \rfloor + 1$ 到 $n$ 之間。
\item \textbf{擅長的操作不同。} 二元搜尋樹可以在 $O(h)$ 時間搜尋任意鍵、找前驅後繼、依序列出；
堆積只能快速拿到最大值（$O(1)$），搜尋任意鍵最差要看過全部 $n$ 個節點。
\end{itemize}

\textbf{㈡ 最小 $n$，最大 $2^n - 1$。} 「以堆積來表示」指用堆積的陣列存法：根放在 $A[1]$，
$A[i]$ 的左子放 $A[2i]$、右子放 $A[2i+1]$（資結書第~7.2.1~節）。 樹缺的位置在陣列裡空著，所以
陣列至少要開到「最大的那個編號」。
\begin{itemize}
\item \textbf{最小值 $n$。} $n$ 個節點至少要 $n$ 格。 當這棵二元搜尋樹剛好是完整二元樹時，編號
恰好是 $1..n$，沒有空洞，開 $n$ 格就夠。 任何 $n$ 都辦得到：先把 $n$ 個位置排成完整二元樹，再照
中序把鍵由小到大填進去，就是一棵二元搜尋樹。
\item \textbf{最大值 $2^n - 1$。} 樹最高是 $n$ 層（每層一個節點），第 $d$ 層（根是第 1 層）的編號
最大是 $2^d - 1$，所以任何節點的編號都不超過 $2^n - 1$。 這個上界達得到：鍵值遞增插入時，每個
節點都是前一個的右子，編號依序是 $1, 3, 7, \dots, 2^n - 1$，陣列要開 $2^n - 1$ 格，只用到其中
$n$ 格。
\end{itemize}
（若陣列從 $A[0]$ 開始存，左右子是 $2i+1$、$2i+2$，答案一樣是 $n$ 和 $2^n - 1$ 格。）

\textbf{㈢ $n + 1$ 個。} 每個節點有 2 個指標欄位，共 $2n$ 個。 除了根以外，每個節點恰好被它的
父節點的一個欄位指著，所以非空的欄位恰好 $n - 1$ 個。 空的欄位有 $2n - (n-1) = n + 1$ 個
（資結書第~7.2.2~節）。

\textbf{㈣ 先建二元搜尋樹。} 依序插入：45 是根；33 $<$ 45 放左；17 放 33 的左；65 放 45 的右；
54 放 65 的左；70 放 65 的右；88 放 70 的右；25：$< 45$、$< 33$、$> 17$，放 17 的右。
中序（由小到大）是 17, 25, 33, 45, 54, 65, 70, 88。

\textbf{再加中序引線}（資結書第~7.5~節）：左指標是空的就改指中序的前一個節點，右指標是空的就改指
中序的下一個節點，並把對應的旗標設為真。
\begin{center}
\begin{tabular}{cll}
\toprule
節點 & 左指標 & 右指標 \\
\midrule
17 & 引線 → \NIL（沒有前驅） & 子節點 25 \\
25 & 引線 → 17 & 引線 → 33 \\
33 & 子節點 17 & 引線 → 45 \\
45 & 子節點 33 & 子節點 65 \\
54 & 引線 → 45 & 引線 → 65 \\
65 & 子節點 54 & 子節點 70 \\
70 & 引線 → 65 & 子節點 88 \\
88 & 引線 → 70 & 引線 → \NIL（沒有後繼） \\
\bottomrule
\end{tabular}
\end{center}
共 9 條引線，正好是㈢ 的 $n + 1 = 9$ 個空指標。 下圖實線是原本的子節點指標，虛線箭頭是引線：
\begin{center}
\begin{tikzpicture}[
    v/.style={circle, draw, minimum size=7.5mm, inner sep=0pt, font=\small},
    th/.style={->, >=stealth, dashed, thick, black!60}]
\node[v] (45) at (0,0) {45};
\node[v] (33) at (-2.6,-1.3) {33};
\node[v] (65) at (2.6,-1.3) {65};
\node[v] (17) at (-4,-2.6) {17};
\node[v] (54) at (1.3,-2.6) {54};
\node[v] (70) at (3.9,-2.6) {70};
\node[v] (25) at (-2.8,-3.9) {25};
\node[v] (88) at (5.1,-3.9) {88};
\draw (45) -- (33); \draw (45) -- (65); \draw (33) -- (17); \draw (17) -- (25);
\draw (65) -- (54); \draw (65) -- (70); \draw (70) -- (88);
\node[font=\scriptsize] (nl) at (-5.4,-3.6) {\NIL};
\node[font=\scriptsize] (nr) at (6.4,-4.9) {\NIL};
\draw[th] (17) -- (nl);
\draw[th] (25) to[out=180, in=-100] (17);
\draw[th] (25) to[out=20, in=-80] (33);
\draw[th] (33) to[out=10, in=-150] (45);
\draw[th] (54) to[out=150, in=-60] (45);
\draw[th] (54) to[out=20, in=-110] (65);
\draw[th] (70) to[out=160, in=-40] (65);
\draw[th] (88) to[out=180, in=-90] (70);
\draw[th] (88) -- (nr);
\end{tikzpicture}
\end{center}
有些課本（例如 Horowitz）會加一個\textbf{頭節點}，讓 17 的左引線和 88 的右引線改指向頭節點，
其餘引線相同。
\end{pastsoln}
